% ============================================================
%  Chapter 3 -- Cauchy Sequences, Upper and Lower Limits,
%               Some Special Sequences: 3.8-3.20
% ============================================================
\rsection{Cauchy Sequences}

\begin{signpost}[Convergence without the limit]
Definition 3.1 cannot be checked without naming the limit in advance --- a
severe handicap, since the whole point of many arguments is to \emph{prove} a
limit exists. This section removes the handicap. A Cauchy sequence is one whose
terms crowd \emph{each other}; Theorem 3.11 says that in $\R^k$ this internal
crowding already guarantees a limit. Rudin says exactly this in the note after
3.11, and it is the single most-used fact in the chapter: every series
convergence test works by verifying the Cauchy condition, never by producing
the sum.
\end{signpost}

\ritem{3.8}{Definition}
A sequence $(p_n)$ in a metric space $X$ is said to be a \emph{Cauchy sequence}
if for every $\varepsilon > 0$ there is an integer $N$ such that
$d(p_n,p_m) < \varepsilon$ if $n \ge N$ and $m \ge N$.\kw{Both indices are
quantified, independently. The weaker condition
$d(p_n,p_{n+1}) \to 0$ --- consecutive terms approaching --- does \emph{not}
suffice: $s_n = \sqrt n$ has $s_{n+1}-s_n \to 0$ yet marches to infinity, and
the partial sums of the harmonic series (3.28) do the same. Mistaking the
consecutive condition for the Cauchy condition is the most common error made
with this definition.}

In our discussion of Cauchy sequences, as well as other situations which will
arise later, the following geometric concept will be useful.

\ritem{3.9}{Definition}
Let $E$ be a nonempty subset of a metric space $X$ and let $S$ be the set of
all real numbers of the form $d(p,q)$ with $p \in E$ and $q \in E$. The sup of
$S$ is called the \emph{diameter} of $E$.

If $(p_n)$ is a sequence in $X$ and if $E_N$ consists of the points
$p_N, p_{N+1}, p_{N+2}, \dots$, it is clear from the two preceding definitions
that $(p_n)$ is a Cauchy sequence if and only if
\[ \lim_{N\to\infty} \operatorname{diam} E_N = 0. \]

\ritem{3.10}{Theorem}
\begin{enumerate}[label=(\alph*), leftmargin=2.2em, itemsep=2pt, topsep=3pt]
  \item \emph{If $\overline{E}$ is the closure of a set $E$ in a metric space
        $X$, then}
        \[ \operatorname{diam} \overline{E} = \operatorname{diam} E. \]
  \item \emph{If $K_n$ is a sequence of compact sets in $X$ such that
        $K_n \supset K_{n+1}$ $(n = 1,2,3,\dots)$ and if}
        \[ \lim_{n\to\infty} \operatorname{diam} K_n = 0, \]
        \emph{then $\bigcap_{n=1}^\infty K_n$ consists of exactly one point.}
\end{enumerate}

\begin{proof}
\begin{enumerate}[label=(\alph*), leftmargin=2.2em, itemsep=4pt, topsep=3pt]
  \item Since $E \subset \overline{E}$, it is clear that
        \[ \operatorname{diam} E \le \operatorname{diam} \overline{E}. \]
        Fix $\varepsilon > 0$, and choose $p \in \overline{E}$,
        $q \in \overline{E}$. By the definition of $\overline{E}$, there are
        points $p'$, $q'$ in $E$ such that $d(p,p') < \varepsilon$,
        $d(q,q') < \varepsilon$. Hence
        \begin{align*}
          d(p,q) &\le d(p,p') + d(p',q') + d(q',q)\\
                 &< 2\varepsilon + d(p',q') \le
                   2\varepsilon + \operatorname{diam} E.
        \end{align*}
        It follows that
        \[ \operatorname{diam} \overline{E} \le 2\varepsilon +
           \operatorname{diam} E, \]
        and since $\varepsilon$ was arbitrary, (a) is proved.
  \item Put $K = \bigcap_{n=1}^\infty K_n$. By Theorem 2.36, $K$ is not empty.
        If $K$ contains more than one point, then $\operatorname{diam} K > 0$.
        But for each $n$, $K_n \supset K$ so that
        $\operatorname{diam} K_n \ge \operatorname{diam} K$. This contradicts
        the assumption that $\operatorname{diam} K_n \to 0$.\kn{Division of
        labour: 2.36 (compactness) supplies \emph{at least} one point;
        shrinking diameters supply \emph{at most} one. Nested-interval
        arguments always have these two halves, and it is worth noticing which
        hypothesis funds which.}
\end{enumerate}
\end{proof}

\ritem{3.11}{Theorem}
\begin{enumerate}[label=(\alph*), leftmargin=2.2em, itemsep=2pt, topsep=3pt]
  \item \emph{In any metric space $X$, every convergent sequence is a Cauchy
        sequence.}
  \item \emph{If $X$ is a compact metric space and if $(p_n)$ is a Cauchy
        sequence in $X$, then $(p_n)$ converges to some point of $X$.}
  \item \emph{In $\R^k$, every Cauchy sequence converges.}
\end{enumerate}

\emph{Note:} The difference between the definition of convergence and the
definition of a Cauchy sequence is that the limit is explicitly involved in the
former, but not in the latter. Thus Theorem 3.11(b) may enable us to decide
whether or not a given sequence converges without knowledge of the limit to
which it may converge.

The fact (contained in Theorem 3.11) that a sequence converges in $\R^k$ if and
only if it is a Cauchy sequence is usually called the \emph{Cauchy criterion}
for convergence.

\begin{proof}
\begin{enumerate}[label=(\alph*), leftmargin=2.2em, itemsep=4pt, topsep=3pt]
  \item If $p_n \to p$ and if $\varepsilon > 0$, there is an integer $N$ such
        that $d(p,p_n) < \varepsilon$ for all $n \ge N$. Hence
        \[ d(p_n,p_m) \le d(p_n,p) + d(p,p_m) < 2\varepsilon \]
        as soon as $n \ge N$ and $m \ge N$. Thus $(p_n)$ is a Cauchy
        sequence.
  \item Let $(p_n)$ be a Cauchy sequence in the compact space $X$. For
        $N = 1,2,3,\dots$, let $E_N$ be the set consisting of
        $p_N, p_{N+1}, p_{N+2}, \dots$. Then
        \begin{equation}
          \lim_{N\to\infty} \operatorname{diam} \overline{E}_N = 0, \tag{3.3}
        \end{equation}
        by Definition 3.9 and Theorem 3.10(a). Being a closed subset of a
        compact space $X$, each $\overline{E}_N$ is compact (Theorem 2.35).
        Also, $E_N \supset E_{N+1}$, so that
        $\overline{E}_N \supset \overline{E}_{N+1}$.

        Theorem 3.10(b) shows now that there is a unique $p \in X$ which lies
        in every $\overline{E}_N$.\kn{Why the closures? $E_N$ itself need not
        be compact --- it is just a scatter of points --- and 3.10(b) needs
        compact sets. Passing to $\overline{E}_N$ costs nothing precisely
        because of 3.10(a): closing a set does not increase its diameter. That
        is the entire reason 3.10(a) was proved.}

        Let $\varepsilon > 0$ be given. By (3.3) there is an integer $N_0$
        such that $\operatorname{diam} \overline{E}_N < \varepsilon$ if
        $N \ge N_0$. Since $p \in \overline{E}_N$, it follows that
        $d(p,q) < \varepsilon$ for every $q \in \overline{E}_N$, hence for
        every $q \in E_N$. In other words, $d(p,p_n) < \varepsilon$ if
        $n \ge N_0$. This says precisely that $p_n \to p$.
  \item Let $(\mathbf{x}_n)$ be a Cauchy sequence in $\R^k$. Define $E_N$ as
        in (b), with $\mathbf{x}_i$ in place of $p_i$. For some $N$,
        $\operatorname{diam} E_N < 1$. The range of $(\mathbf{x}_n)$ is the
        union of $E_N$ and the finite set
        $\{\mathbf{x}_1,\dots,\mathbf{x}_{N-1}\}$. Hence $(\mathbf{x}_n)$ is
        bounded. Since every bounded subset of $\R^k$ has compact closure in
        $\R^k$ (Theorem 2.41), (c) follows from (b).\kn{The structure of (c):
        Cauchy $\Rightarrow$ bounded (cheaply, from one $E_N$ of finite
        diameter plus finitely many strays --- the same split as 3.2(c)), then
        bounded $\Rightarrow$ inside a compact set (2.41), then apply (b). All
        the completeness is imported through 2.41.}
\end{enumerate}
\end{proof}

\begin{center}
\begin{tikzpicture}[x=1mm, y=1mm]
  \node[rmbox, text width=26mm] (b0) at (0,22)
    {$(p_n)$ Cauchy};
  \node[rmbox, text width=30mm] (b1) at (43,22)
    {$\operatorname{diam} E_N \to 0$\\ \textit{(3.9, tails shrink)}};
  \node[rmbox, text width=30mm] (b2) at (89,22)
    {$\overline{E}_N$: compact,\\nested, diam $\to 0$};
  \draw[rmarrow] (b0) -- (b1);
  \draw[rmarrow] (b1) -- node[lbl, anchor=south, inner sep=2.5pt]
    {2.35, 3.10(a)} (b2);
  \node[rmbox, text width=30mm] (b3) at (89,0)
    {one $p$ in every $\overline{E}_N$};
  \draw[rmarrow] (b2) -- node[lbl, anchor=west, inner sep=3pt]
    {3.10(b)} (b3);
  \node[rmbox, text width=33mm] (b4) at (40,0)
    {$d(p,p_n) \le \operatorname{diam}\overline{E}_N$:\\$p_n \to p$};
  \draw[rmarrow] (b3) -- (b4);
\end{tikzpicture}
\end{center}
\figcap{Theorem 3.11(b) as a flow diagram: the tail-sets carry the whole
argument, closures make them compact, and the nested-set theorem hands over
the limit. Part (c) bolts ``Cauchy $\Rightarrow$ bounded $\Rightarrow$ inside
a compact set (2.41)'' onto the front.}

\ritem{3.12}{Definition}
A metric space in which every Cauchy sequence converges is said to be
\emph{complete}.

Thus \emph{Theorem 3.11 says that all compact metric spaces and all euclidean
spaces are complete}. Theorem 3.11 implies also that \emph{every closed subset
$E$ of a complete metric space $X$ is complete}. (Every Cauchy sequence in $E$
is a Cauchy sequence in $X$, hence it converges to some $p \in X$, and actually
$p \in E$ since $E$ is closed.) An example of a metric space which is not
complete is the space of all rational numbers, with $d(x,y) = |x-y|$.

\begin{deeper}[The word ``complete,'' finally --- and the full circle]
This is the definition promised since Chapter 1. The comparison box at 1.10
said the least-upper-bound property ``is what the word \emph{complete} will
name in 3.12''; here is the naming, and with it the last of the standard
completeness statements enters the book. Of the ECNU text's six equivalent
formulations, five have now appeared: the supremum principle (1.19), nested
intervals (2.38), Heine--Borel (2.41), Bolzano--Weierstrass (2.42 and 3.6(b)),
and now the Cauchy criterion (3.11(c)); monotone convergence follows as 3.14.

There is also a genuine full circle here. Tao \emph{constructs} $\R$ as the
completion of $\Q$ --- real numbers are, in his book, equivalence classes of
Cauchy sequences of rationals, so ``every Cauchy sequence converges'' is nearly
true by construction. Rudin travelled the other way: $\R$ was characterised by
suprema, and Cauchy completeness arrived at the end of a chain through
compactness. Exercise 3.24 closes the loop from Rudin's side, carrying out the
Cauchy-sequence completion for an arbitrary metric space.

One caution: completeness and compactness part company here. Compact implies
complete (3.11(b)), but $\R^k$ is complete without being compact, and
completeness is \emph{not} intrinsic in the sense of 2.33 --- $(0,1)$ and $\R$
are homeomorphic, yet one is complete and the other not. Completeness belongs
to the metric, not the topology.
\end{deeper}

Theorem 3.2(c) and example (d) after Definition 3.1 show that convergent
sequences are bounded, but that bounded sequences in $\R^k$ need not converge.
However, there is one important case in which convergence is equivalent to
boundedness; this happens for monotonic sequences in $\R^1$.

\ritem{3.13}{Definition}
A sequence $(s_n)$ of real numbers is said to be
\begin{enumerate}[label=(\alph*), leftmargin=2.2em, itemsep=1pt, topsep=3pt]
  \item \emph{monotonically increasing} if $s_n \le s_{n+1}$
        $(n = 1,2,3,\dots)$;
  \item \emph{monotonically decreasing} if $s_n \ge s_{n+1}$
        $(n = 1,2,3,\dots)$.
\end{enumerate}
The class of \emph{monotonic} sequences consists of the increasing and the
decreasing sequences.

\ritem{3.14}{Theorem}
\emph{Suppose $(s_n)$ is monotonic. Then $(s_n)$ converges if and only if it is
bounded.}

\begin{proof}
Suppose $s_n \le s_{n+1}$ (the proof is analogous in the other case). Let $E$
be the range of $(s_n)$. If $(s_n)$ is bounded, let $s$ be the least upper
bound of $E$. Then
\[ s_n \le s \qquad (n = 1,2,3,\dots). \]
For every $\varepsilon > 0$, there is an integer $N$ such that
\[ s - \varepsilon < s_N \le s, \]
for otherwise $s - \varepsilon$ would be an upper bound of $E$.\kn{Clause (ii)
of Definition 1.8 yet again, in its working form --- step back from the
supremum and some term sticks out past you. Monotonicity then does what no
other hypothesis could: one term beyond $s-\varepsilon$ drags the entire tail
with it.} Since $(s_n)$ increases, $n \ge N$ therefore implies
\[ s - \varepsilon < s_n \le s, \]
which shows that $(s_n)$ converges (to $s$).

The converse follows from Theorem 3.2(c).
\end{proof}

\begin{center}
\begin{tikzpicture}[x=1mm, y=1mm]
  \draw[axis] (-2,0) -- (96,0);
  \node[lbl, anchor=north] at (94,-2) {$n$};
  % sup line
  \draw[thin] (0,22) -- (92,22);
  \node[mlbl, anchor=east] at (-2,22) {$s=\sup E$};
  \draw[guide] (0,17) -- (92,17);
  \node[mlbl, anchor=west] at (92.5,17) {$s-\varepsilon$};
  % staircase terms
  \foreach \x/\y in {5/4, 13/8, 21/11, 29/13.5, 37/15.5, 45/18, 53/19.3,
                     61/20.2, 69/20.8, 77/21.2, 85/21.5}{
    \node[ptfill, minimum size=2.2pt] at (\x,\y) {};}
  \draw[guide] (45,-4.5) -- (45,18);
  \node[lbl, anchor=north] at (45,-5.5) {$N$: first term past $s-\varepsilon$};
  \node[lbl, anchor=south] at (26,24.5) {increasing, never past $s$};
\end{tikzpicture}
\end{center}
\figcap{Monotone convergence: the terms climb, the supremum caps them, and once
one term passes $s-\varepsilon$ the rest are trapped in the gap. The limit
\emph{is} the supremum --- 3.14 is Definition 1.10 restated for sequences,
which is why the ECNU text lists the two as equivalent completeness axioms.}

\rsection{Upper and Lower Limits}

\begin{signpost}[The hardest definition in the chapter, and what it is for]
A divergent sequence can still have structure: $(-1)^n$ oscillates between two
values rather than wandering. $\limsup$ and $\liminf$ capture the outer
envelope of that oscillation --- the largest and smallest subsequential
limits. Their virtue is \emph{unconditional existence}: every real sequence has
both, in the extended reals, with no hypothesis at all. That is why the root
test (3.33) and the radius-of-convergence formula (3.39) are stated with
$\limsup$ --- an ordinary limit might not exist, and the theory would fill up
with exceptions.

\smallskip
Keep Theorem 3.17 in view as the working characterisation: $s^* $ is the one
number that is both \emph{reached} (along some subsequence) and \emph{barely
exceeded} (only finitely many terms above any $x > s^*$). The definition via
$\sup E$ is for existence; 3.17 is what proofs actually use.
\end{signpost}

\ritem{3.15}{Definition}
Let $(s_n)$ be a sequence of real numbers with the following property: For
every real $M$ there is an integer $N$ such that $n \ge N$ implies
$s_n \ge M$. We then write
\[ s_n \to +\infty. \]
Similarly, if for every real $M$ there is an integer $N$ such that $n \ge N$
implies $s_n \le M$, we write
\[ s_n \to -\infty. \]

It should be noted that we now use the symbol $\to$ (introduced in Definition
3.1) for certain types of divergent sequences, as well as for convergent
sequences, but that the definitions of convergence and of limit, given in
Definition 3.1, are unchanged.\kw{So ``$s_n \to +\infty$'' and ``$(s_n)$
diverges'' are compatible statements --- the arrow has been extended, the word
\emph{limit} has not. This is the extended-real bookkeeping of 1.23 paying its
way: $+\infty$ is a destination for notation, not a number a sequence can
converge to.}

\ritem{3.16}{Definition}
Let $(s_n)$ be a sequence of real numbers. Let $E$ be the set of numbers $x$
(in the extended real number system) such that $s_{n_k} \to x$ for some
subsequence $(s_{n_k})$. This set $E$ contains all subsequential limits as
defined in Definition 3.5, plus possibly the numbers $+\infty$, $-\infty$.

We now recall Definitions 1.8 and 1.23 and put
\[ s^* = \sup E, \qquad s_* = \inf E. \]
The numbers $s^*$, $s_*$ are called the \emph{upper} and \emph{lower limits} of
$(s_n)$; we use the notation
\[ \limsup_{n\to\infty} s_n = s^*, \qquad \liminf_{n\to\infty} s_n = s_*. \]

\begin{center}
\begin{tikzpicture}[x=1mm, y=1mm]
  \draw[axis] (-2,0) -- (98,0);
  \node[lbl, anchor=north] at (96,-2) {$n$};
  % envelopes
  \draw[thin] (0,24) -- (94,24);
  \node[mlbl, anchor=west] at (94.5,24) {$s^*=1$};
  \draw[thin] (0,4) -- (94,4);
  \node[mlbl, anchor=west] at (94.5,4) {$s_*=-1$};
  % oscillating terms approaching the envelopes: s_n=(-1)^n/(1+1/n) scaled
  % top subsequence rises to 24, bottom falls to 4  (centre line at 14)
  \foreach \x/\y in {6/19.0, 18/21.6, 30/22.6, 42/23.0, 54/23.3, 66/23.5, 78/23.6, 90/23.7}{
    \node[ptfill, minimum size=2.2pt] at (\x,\y) {};}
  \foreach \x/\y in {12/7.7, 24/5.7, 36/5.0, 48/4.7, 60/4.5, 72/4.4, 84/4.35}{
    \node[ptopen, minimum size=2.6pt] at (\x,\y) {};}
  \node[lbl, anchor=south] at (30,27) {even terms $\to 1$};
  \node[lbl, anchor=north] at (30,-3) {odd terms $\to -1$};
\end{tikzpicture}
\end{center}
\figcap{Example 3.18(b), $s_n = (-1)^n/(1+1/n)$: two subsequences, two
subsequential limits, $E = \{-1,1\}$. The upper and lower limits are the
envelope the oscillation settles into --- not the largest and smallest
\emph{terms}.}

\ritem{3.17}{Theorem}
\emph{Let $(s_n)$ be a sequence of real numbers. Let $E$ and $s^*$ have the
same meaning as in Definition 3.16. Then $s^*$ has the following two
properties:}
\begin{enumerate}[label=(\alph*), leftmargin=2.2em, itemsep=1pt, topsep=3pt]
  \item \emph{$s^* \in E$.}
  \item \emph{If $x > s^*$, there is an integer $N$ such that $n \ge N$ implies
        $s_n < x$.}
\end{enumerate}
\emph{Moreover, $s^*$ is the only number with the properties (a) and (b).}

\emph{Of course, an analogous result is true for $s_*$.}

\begin{proof}
\begin{enumerate}[label=(\alph*), leftmargin=2.2em, itemsep=4pt, topsep=3pt]
  \item If $s^* = +\infty$, then $E$ is not bounded above; hence $(s_n)$ is
        not bounded above, and there is a subsequence $(s_{n_k})$ such that
        $s_{n_k} \to +\infty$.

        If $s^*$ is real, then $E$ is bounded above, and at least one
        subsequential limit exists, so that (a) follows from Theorems 3.7 and
        2.28.\kn{Here is where the section's stockpile is spent, both pieces
        at once: 3.7 says $E$ is closed, 2.28 says $\sup E$ lies in
        $\overline{E} = E$. So the supremum of the subsequential limits is
        itself a subsequential limit --- attained, not merely approached. This
        is the fact that makes $\limsup$ usable.}

        If $s^* = -\infty$, then $E$ contains only one element, namely
        $-\infty$, and there is no subsequential limit. Hence, for any real
        $M$, $s_n > M$ for at most a finite number of values of $n$, so that
        $s_n \to -\infty$.

        This establishes (a) in all cases.
  \item Suppose there is a number $x > s^*$ such that $s_n \ge x$ for
        infinitely many values of $n$. In that case, there is a number
        $y \in E$ such that $y \ge x > s^*$, contradicting the definition of
        $s^*$.\kn{The compressed step: infinitely many terms $\ge x$ form a
        subsequence in $[x,+\infty]$, and 3.6/3.17(a) applied to it yields a
        subsequential limit $y \ge x$. ``Infinitely many terms above $x$''
        always hands you an element of $E$ at or above $x$ --- that
        conversion, bounded-infinite-to-limit, is Bolzano--Weierstrass 3.6(b)
        on duty.}
\end{enumerate}
Thus $s^*$ satisfies (a) and (b). To show the uniqueness, suppose there are
two numbers, $p$ and $q$, which satisfy (a) and (b), and suppose $p < q$.
Choose $x$ such that $p < x < q$. Since $p$ satisfies (b), we have $s_n < x$
for $n \ge N$. But then $q$ cannot satisfy (a).
\end{proof}

\begin{unpack}[Three working forms of one number]
$\limsup$ wears three interchangeable outfits, and fluency means switching
freely.

\smallskip
\textbf{As $\sup E$} (Definition 3.16): best for existence and for symmetry
with $\liminf$.

\smallskip
\textbf{As the two-sided test} (Theorem 3.17): $s^*$ is reached along a
subsequence, and any $x > s^*$ eventually caps the whole sequence. In
$\varepsilon$-form: for every $\varepsilon > 0$, $s_n > s^*+\varepsilon$ only
finitely often, while $s_n > s^*-\varepsilon$ infinitely often. This is the
form the root test uses.

\smallskip
\textbf{As a limit of suprema}: $s^* = \lim_{N\to\infty} \sup_{n\ge N} s_n$
--- the tail-suprema $\sup_{n \ge N} s_n$ decrease as the tail shrinks, so
they converge by 3.14, and their limit is $s^*$. Rudin leaves this form to
Exercise 3.5's neighbourhood; Tao takes it as the \emph{definition} (his
6.4.6, via ``$\sup$ of the tail''), and the ECNU text develops it in its
\S7.2. It explains the name: the \emph{lim} of the \emph{sup}s.
\end{unpack}

\ritem{3.18}{Examples}
\begin{enumerate}[label=(\alph*), leftmargin=2.2em, itemsep=3pt, topsep=4pt]
  \item Let $(s_n)$ be a sequence containing all rationals. Then every real
        number is a subsequential limit, and
        \[ \limsup_{n\to\infty} s_n = +\infty, \qquad
           \liminf_{n\to\infty} s_n = -\infty. \]\kn{Such a sequence exists
        because $\Q$ is countable (Corollary to 2.13), and every real is a
        subsequential limit because $\Q$ is dense (1.20(b)) --- the two
        headline facts about $\Q$, cooperating. Here $E$ is all of the
        extended line: divergence as total as it gets.}
  \item Let $s_n = (-1)^n/[1+(1/n)]$. Then
        \[ \limsup_{n\to\infty} s_n = 1, \qquad
           \liminf_{n\to\infty} s_n = -1. \]
  \item For a real-valued sequence $(s_n)$, $\lim_{n\to\infty} s_n = s$ if and
        only if
        \[ \limsup_{n\to\infty} s_n = \liminf_{n\to\infty} s_n = s. \]\kn{The
        squeeze that makes the pair useful: convergence is exactly ``the
        envelope closes.'' Many later proofs show $\limsup \le c$ and
        $\liminf \ge c$ separately and conclude $\lim = c$; 3.19 exists to
        service such arguments.}
\end{enumerate}

We close this section with a theorem which is useful and whose proof is quite
trivial:

\ritem{3.19}{Theorem}
\emph{If $s_n \le t_n$ for $n \ge N$, where $N$ is fixed, then}
\[ \liminf_{n\to\infty} s_n \le \liminf_{n\to\infty} t_n, \qquad
   \limsup_{n\to\infty} s_n \le \limsup_{n\to\infty} t_n. \]

\rsection{Some Special Sequences}

\begin{signpost}[A toolkit, assembled by one method]
Five limits used from here to the end of the book, each proved by the squeeze
stated in the opening line: trap $x_n$ between $0$ and something already known
to vanish. The recurring supplier of the trap is the binomial theorem ---
expand $(1+x)^n$, keep a single term, and discard the rest. Item (d) is the
one to remember: powers beat polynomials, and (e), geometric decay, is its
special case $\alpha = 0$. These feed directly into 3.26 and every comparison
with a geometric series afterwards.
\end{signpost}

\noindent
We shall now compute the limits of some sequences which occur frequently. The
proofs will all be based on the following remark: \emph{If $0 \le x_n \le s_n$
for $n \ge N$, where $N$ is some fixed number, and if $s_n \to 0$, then
$x_n \to 0$.}

\ritem{3.20}{Theorem}
\begin{enumerate}[label=(\alph*), leftmargin=2.2em, itemsep=2pt, topsep=3pt]
  \item \emph{If $p > 0$, then $\displaystyle\lim_{n\to\infty} \frac{1}{n^p} = 0$.}
  \item \emph{If $p > 0$, then $\displaystyle\lim_{n\to\infty} \sqrt[n]{p} = 1$.}
  \item \emph{$\displaystyle\lim_{n\to\infty} \sqrt[n]{n} = 1$.}
  \item \emph{If $p > 0$ and $\alpha$ is real, then
        $\displaystyle\lim_{n\to\infty} \frac{n^\alpha}{(1+p)^n} = 0$.}
  \item \emph{If $|x| < 1$, then $\displaystyle\lim_{n\to\infty} x^n = 0$.}
\end{enumerate}

\begin{proof}
\begin{enumerate}[label=(\alph*), leftmargin=2.2em, itemsep=4pt, topsep=3pt]
  \item Take $n > (1/\varepsilon)^{1/p}$. (Note that the archimedean property
        of the real number system is used here.)\kn{And note \emph{what else}
        is used: $n^p$ for real $p$ is Exercise 1.6's construction, and the
        inequality manipulation needs 1.21. The one-line proof stands on the
        chapter-1 exercises.}
  \item If $p > 1$, put $x_n = \sqrt[n]{p} - 1$. Then $x_n > 0$ and, by the
        binomial theorem,
        \[ 1 + nx_n \le (1+x_n)^n = p, \]
        so that
        \[ 0 < x_n \le \frac{p-1}{n}. \]
        Hence $x_n \to 0$. If $p = 1$, (b) is trivial, and if $0 < p < 1$, the
        result is obtained by taking reciprocals.
  \item Put $x_n = \sqrt[n]{n} - 1$. Then $x_n \ge 0$ and, by the binomial
        theorem,
        \[ n = (1+x_n)^n \ge \frac{n(n-1)}{2}x_n^2. \]
        Hence
        \[ 0 \le x_n \le \sqrt{\frac{2}{n-1}} \qquad (n \ge 2). \]\kn{Same
        trick as (b), one term deeper. Keeping the linear term gives only
        $x_n \le (n-1)/n$, which is useless; the quadratic term grows like
        $n^2$ and beats the left side. Choosing \emph{which} binomial term to
        keep is the entire skill, and (d) makes the choice a third time, with
        term $k$.}
  \item Let $k$ be an integer such that $k > \alpha$, $k > 0$. For $n > 2k$,
        \[ (1+p)^n > \binom{n}{k}p^k =
           \frac{n(n-1)\cdots(n-k+1)}{k!}\,p^k > \frac{n^kp^k}{2^kk!}. \]
        Hence
        \[ 0 < \frac{n^\alpha}{(1+p)^n} < \frac{2^kk!}{p^k}\,n^{\alpha-k}
           \qquad (n > 2k). \]
        Since $\alpha - k < 0$, $n^{\alpha-k} \to 0$, by (a).
  \item Take $\alpha = 0$ in (d).
\end{enumerate}
\end{proof}
